% Part: turing-machines
% Chapter: machines-computations
% Section: representing-tms

\documentclass[../../../include/open-logic-section]{subfiles}

\begin{document}

\olfileid{tur}{mac}{rep}
\olsection{টুরিং যন্ত্রের উপস্থাপন}

\begin{explain}
টুরিং যন্ত্রকে \emph{দশা-চিত্র} দিয়ে দৃশ্যত উপস্থাপন করা যায়। চিত্রগুলি
তীর দিয়ে যুক্ত দশা-নোড দিয়ে গঠিত। স্বাভাবিকভাবেই, প্রতিটি দশা-নোড
যন্ত্রের একটি দশা প্রকাশ করে। প্রতিটি তীর সেই দশা থেকে পালনযোগ্য একটি
নির্দেশ প্রকাশ করে এবং নির্দেশটির বিবরণ সংশ্লিষ্ট তীরের উপরে বা নিচে
লেখা থাকে। নিচের যন্ত্রটি বিবেচনা করো; এর মাত্র দুটি অভ্যন্তরীণ দশা,
$q_0$ ও $q_1$, এবং একটি নির্দেশ আছে:
\[
\begin{tikzpicture}[->,>=stealth',shorten >=1pt,auto,node distance=2.8cm,
                    semithick]
  \tikzstyle{every state}=[fill=none,draw=black,text=black]

  \node[initial,state]   (A)              {$q_0$};
  \node[state]   (B) [right of=A] {$q_1$};

  \path (A) edge  node {\TMtrans{\TMblank}{\TMstroke}{\TMright}} (B);
\end{tikzpicture}
\]
মনে রেখো, টুরিং যন্ত্রের একটি পাঠ-লেখ হেড এবং নিবেশ-লেখা একটি ফিতা
আছে। নির্দেশটি এভাবে পড়া যায়: \emph{দশা~$q_0$-তে একটি~$\TMblank$
পড়লে একটি~$\TMstroke$ লেখো, ডান দিকে যাও এবং দশা~$q_1$-এ যাও}। এটি
অবস্থান্তর অপেক্ষক দ্বারা $\tuple{q_0, \TMblank}$-কে
$\tuple{q_1, \TMstroke, \TMright}$-এ পাঠানোর সমতুল্য।
\end{explain}

\begin{ex}
\emph{জোড়-যন্ত্র}: নিচের টুরিং যন্ত্রটি ঠিক তখনই থামে, যখন ফিতায়
$\TMstroke$-এর সংখ্যা জোড় (ধরা হচ্ছে, ফিতার প্রথম $\TMblank$-এর আগে
সব $\TMstroke$ রয়েছে)।
\[
\begin{tikzpicture}[->,>=stealth',shorten >=1pt,auto,node distance=2.8cm,
                    semithick]
  \tikzstyle{every state}=[fill=none,draw=black,text=black]

  \node[initial,state]         (A)              {$q_0$};
  \node[state]         (B) [right of=A] {$q_1$};

  \path (A) edge [bend left] node {\TMtrans{\TMstroke}{\TMstroke}{\TMright}} (B)
        (B) edge [loop above] node {\TMtrans{\TMblank}{\TMblank}{\TMright}} (B)
            edge [bend left] node {\TMtrans{\TMstroke}{\TMstroke}{\TMright}} (A);
\end{tikzpicture}
\]

দশা-চিত্রটি নিচের অবস্থান্তর অপেক্ষকের সঙ্গে সঙ্গতিপূর্ণ:
\begin{align*}
\delta(q_0, \TMstroke) & = \tuple{q_1, \TMstroke, \TMright},\\
\delta(q_1, \TMstroke) & = \tuple{q_0, \TMstroke, \TMright},\\
\delta(q_1, \TMblank)  & = \tuple{q_1, \TMblank, \TMright}
\end{align*}
\end{ex}

\begin{explain}
উপরের যন্ত্রটি শুধু তখনই থামে, যখন নিবেশটি জোড়সংখ্যক দাগ। অন্যথায়
যন্ত্রটি (তত্ত্বগতভাবে) অনির্দিষ্টকাল চলতে থাকে। যেকোনো যন্ত্র ও নিবেশের
ক্ষেত্রে যন্ত্রটির \emph{কনফিগারেশন}-গুলি ধাপে ধাপে অনুসরণ করে নির্গম
নির্ণয় করা যায়। পরে আমরা কনফিগারেশনের আনুষ্ঠানিক সংজ্ঞা দেব। আপাতত
স্বজ্ঞাতভাবে কনফিগারেশনকে এমন কয়েকটি চিত্রের অনুক্রম ভাবতে পারি, যা
চালনার যেকোনো মুহূর্তে যন্ত্রের দশা দেখায়। কনফিগারেশনে ফিতার বিষয়বস্তু,
যন্ত্রের দশা এবং পাঠ-লেখ হেডের অবস্থান দেখা যায়।

জোড়-যন্ত্রটিকে চারটি $\TMstroke$-এর নিবেশ দিয়ে চালু করলে তার
কনফিগারেশনগুলি অনুসরণ করি। এক্ষেত্রে যন্ত্রটি থামবে বলে আমাদের প্রত্যাশা।
তারপর যন্ত্রটিকে তিনটি $\TMstroke$-এর নিবেশে চালাব; সেখানে যন্ত্রটি
চিরকাল চলবে।

যন্ত্রটি দশা~$q_0$-তে বাঁদিকের প্রথম~$\TMstroke$ পড়তে পড়তে শুরু করে।
যন্ত্রটির আরম্ভিক দশা নিচের মতো প্রকাশ করা যায়:
\[
\TMendtape \TMstroke_0 \TMstroke \TMstroke \TMstroke \TMblank \ldots
\]
উপরের কনফিগারেশনটি সরল। দেখা যাচ্ছে, যন্ত্রটি দশা~$q_0$-তে বাঁদিকের
প্রথম~$\TMstroke$ পড়তে পড়তে শুরু করে। প্রথম~$\TMstroke$-এ দশার নামের
নিম্নসূচক দিয়ে তা দেখানো হয়েছে। এই মুহূর্তে প্রযোজ্য নির্দেশটি হলো
$\delta(q_0, \TMstroke) = \tuple{q_1, \TMstroke, \TMright}$; ফলে
যন্ত্রটি ফিতায় ডান দিকে যায় এবং দশা~$q_1$-এ বদলে যায়।
% BN-SRC-163 TARGET-CORRECTION-BEGIN
\par\smallskip\noindent\textnormal{\small\textbf{উৎস-সংশোধন (BN-SRC-163):} হিমায়িত ইংরেজি বর্ণনায় এই আরম্ভিক কনফিগারেশনকে “state one” বলা হয়েছে, যদিও আগের বাক্য, চিত্রের নিম্নসূচক এবং প্রযোজ্য নির্দেশ—সবই দশা~$q_0$ নির্দিষ্ট করে। বাংলা পাঠে দশা~$q_0$ রাখা হয়েছে।} % BN-SRC-163 TARGET-CORRECTION-END
\[
\TMendtape \TMstroke \TMstroke_1 \TMstroke \TMstroke \TMblank \ldots
\]
যন্ত্রটি এখন দশা~$q_1$-এ একটি~$\TMstroke$ পড়ছে বলে আমাদের নির্দেশ
$\delta(q_1, \TMstroke) = \tuple{q_0, \TMstroke, \TMright}$ “অনুসরণ”
করতে হবে। এর ফলে কনফিগারেশনটি হয়
\[
\TMendtape \TMstroke \TMstroke \TMstroke_0 \TMstroke \TMblank \ldots
\]
যন্ত্রটি চলতে থাকলে একই ক্রমে নিয়মগুলি আবার প্রয়োগ হয় এবং নিচের দুটি
কনফিগারেশন পাওয়া যায়:
\[
\TMendtape \TMstroke \TMstroke \TMstroke \TMstroke_1 \TMblank \ldots
\]
\[
\TMendtape \TMstroke \TMstroke \TMstroke \TMstroke \TMblank_0 \ldots
\]
যন্ত্রটি এখন দশা~$q_0$-তে একটি~$\TMblank$ পড়ছে। অবস্থান্তর-চিত্র থেকে
সহজেই দেখা যায়, পালন করার কোনো নির্দেশ নেই; তাই যন্ত্রটি থেমে গেছে। এর
অর্থ, নিবেশটি গৃহীত হয়েছে।

এবার ধরা যাক, যন্ত্রটিকে তিনটি $\TMstroke$-এর নিবেশ দিয়ে চালু করি। একই
নির্দেশগুলি পালিত হয় বলে প্রথম কয়েকটি কনফিগারেশন একই রকম; কেবল ফিতার
নিবেশে সামান্য পার্থক্য:
\[
\TMendtape \TMstroke_0 \TMstroke \TMstroke \TMblank \ldots
\]
\[
\TMendtape \TMstroke \TMstroke_1 \TMstroke \TMblank \ldots
\]
\[
\TMendtape \TMstroke \TMstroke \TMstroke_0 \TMblank \ldots
\]
\[
\TMendtape \TMstroke \TMstroke \TMstroke \TMblank_1 \ldots
\]
যন্ত্রটি এখন সব $\TMstroke$ পেরিয়ে গেছে এবং দশা~$q_1$-এ একটি
~$\TMblank$ পড়ছে। চিত্রে দেখানো আছে, এক্ষেত্রে
$\delta(q_1, \TMblank) =\tuple{q_1, \TMblank, \TMright}$ আকারের একটি
নির্দেশ রয়েছে। ফিতাটি ডান দিকে অনির্দিষ্টকাল $\TMblank$ দিয়ে ভরা বলে
যন্ত্রটি দশা~$q_1$-এ থেকে এবং ক্রমশ আরও ডান দিকে যেতে যেতে এই নির্দেশটি
\emph{চিরকাল} পালন করবে। যন্ত্রটি কখনোই থামবে না এবং নিবেশটি গ্রহণ করবে
না।
\end{explain}

\begin{explain}
মনে রাখা জরুরি যে সব যন্ত্র থামে না। থামা যদি যন্ত্রটির পালনযোগ্য নির্দেশ
ফুরিয়ে যাওয়া বোঝায়, তবে প্রতিটি দশায় প্রতিটি প্রতীকের জন্য একটি করে
বহির্মুখী তীর নিশ্চিত করলেই কখনও না-থামা যন্ত্র বানানো যায়। দশা~$q_0$-তে
একটি $\TMblank$ পড়ার নির্দেশ যোগ করে জোড়-যন্ত্রটিকে অনির্দিষ্টকাল চলার
মতো বদলে নেওয়া যায়।
\end{explain}

\begin{ex}
\[
\begin{tikzpicture}[->,>=stealth',shorten >=1pt,auto,node distance=2.8cm,
                    semithick]
  \tikzstyle{every state}=[fill=none,draw=black,text=black]

  \node[initial,state] (A)              {$q_0$};
  \node[state]         (B) [right of=A] {$q_1$};

  \path
  (A) edge [bend left] node {\TMtrans{\TMstroke}{\TMstroke}{\TMright}} (B)
      edge [loop above] node {\TMtrans{\TMblank}{\TMblank}{\TMright}} (A)
  (B) edge [loop above] node {\TMtrans{\TMblank}{\TMblank}{\TMright}} (B)
      edge [bend left] node {\TMtrans{\TMstroke}{\TMstroke}{\TMright}} (A);
\end{tikzpicture}
\]
\end{ex}

\begin{explain}
টুরিং যন্ত্র উপস্থাপনের আরেকটি উপায় হলো যন্ত্র-সারণি। যন্ত্র-সারণির
$x$-অক্ষে ফিতার বর্ণমালা এবং $y$-অক্ষে যন্ত্রের দশাগুলির সেট দেখানো হয়।
সারণিতে প্রতিটি দশা ও প্রতীকের ছেদস্থলে নির্দেশটির বাকি অংশ—নতুন দশা,
নতুন প্রতীক এবং চলনের দিক—লেখা হয়। কোন দশায় এবং কোন প্রতীকের জন্য
যন্ত্রটি থামে, তা যন্ত্র-সারণি থেকে সহজে নির্ণয় করা যায়। সারণির প্রতিটি
ফাঁকই যন্ত্রটির থামার একটি সম্ভাব্য স্থান। দশা-চিত্র ও নির্দেশ-সেটে
যন্ত্রটির থামার স্থানগুলি সবসময় সঙ্গে সঙ্গে চোখে না পড়লেও, যন্ত্র-সারণির
ফাঁকগুলি খুঁজে যেকোনো থামার স্থান দ্রুত শনাক্ত করা যায়।
\end{explain}

\begin{ex}
জোড়-যন্ত্রটির যন্ত্র-সারণি হলো:
\[
\centering
\begin{tabular}{lllll}
\cline{2-4}
\multicolumn{1}{l|}{}      & \multicolumn{1}{c|}{$\TMblank$}
& \multicolumn{1}{c|}{$\TMstroke$}     & \multicolumn{1}{c|}{$\TMendtape$}   &  \\ \cline{1-4}
\multicolumn{1}{|l|}{$q_0$} & \multicolumn{1}{l|}{}
& \multicolumn{1}{l|}{\TMtrans{q_1}{\TMstroke}{\TMright}} &
\multicolumn{1}{l|}{\phantom{\TMtrans{\TMstroke}{q_1}{\TMright}}} &  \\ \cline{1-4}
\multicolumn{1}{|l|}{$q_1$} & \multicolumn{1}{l|}{\TMtrans{q_1}{\TMblank}{\TMright}}
& \multicolumn{1}{l|}{\TMtrans{q_0}{\TMstroke}{\TMright}} &
\multicolumn{1}{l|}{\phantom{\TMtrans{\TMstroke}{q_1}{\TMright}}} &  \\ \cline{1-4}
\end{tabular}
\]
দেখা যাচ্ছে, দশা~$q_0$-তে একটি~$\TMblank$ পড়ার সময় যন্ত্রটি থামে।
% BN-SRC-164 TARGET-CORRECTION-BEGIN
\par\smallskip\noindent\textnormal{\small\textbf{উৎস-সংশোধন (BN-SRC-164):} হিমায়িত ইংরেজি সারণির বর্ণনা অনুযায়ী প্রতিটি ঘরে নতুন দশা, নতুন প্রতীক ও চলনের দিক—এই ক্রমে লেখা হওয়ার কথা; কিন্তু সক্রিয় তিনটি ঘরেই নতুন প্রতীক ও নতুন দশা উল্টো ক্রমে আছে। বাংলা সারণিতে যথাক্রমে $\tuple{q_1,\TMstroke,\TMright}$, $\tuple{q_1,\TMblank,\TMright}$ ও $\tuple{q_0,\TMstroke,\TMright}$ রাখা হয়েছে, যা আগের অবস্থান্তর অপেক্ষকের সঙ্গে মেলে।} % BN-SRC-164 TARGET-CORRECTION-END
\end{ex}

\begin{explain}
এ পর্যন্ত আমরা কেবল নিবেশ পড়ে গ্রহণ করে এমন যন্ত্র বিবেচনা করেছি। কিন্তু
টুরিং যন্ত্র পড়া ও লেখা দুটিই করতে পারে। এমন যন্ত্রের একটি উদাহরণ (যদিও
উদাহরণের অভাব নেই) হলো \emph{দ্বিগুণকারী}। ফিতায় $n$টি
$\TMstroke$-এর একটি খণ্ড দিয়ে দ্বিগুণকারীকে চালু করলে সেটি $2n$টি
$\TMstroke$-এর একটি খণ্ড নির্গত করে।
\end{explain}

\begin{ex}
  \ollabel{ex:doubler} দ্বিগুণকারী যন্ত্র নির্মাণের আগে সমস্যাটি সমাধানের
  একটি \emph{কৌশল} ঠিক করা জরুরি। যন্ত্রটি (আমরা যেভাবে সূত্রবদ্ধ করেছি)
  কতগুলি $\TMstroke$ পড়েছে তা মনে রাখতে পারে না; তাই ফিতার সব
  $\TMstroke$-এর হিসাব রাখার একটি উপায় দরকার। তার একটি উপায় হলো একটি
  ~$\TMblank$ দিয়ে নির্গমকে নিবেশ থেকে আলাদা করা। যন্ত্রটি তখন নিবেশের
  প্রথম $\TMstroke$ মুছতে পারে, নিবেশের বাকি অংশ পেরিয়ে যেতে পারে,
  একটি $\TMblank$ রেখে দুটি নতুন $\TMstroke$ লিখতে পারে। তারপর যন্ত্রটি
  ফিরে গিয়ে নিবেশের দ্বিতীয় $\TMstroke$ খুঁজে সেটিকেও দ্বিগুণ করবে।
  নিবেশের প্রতিটি $\TMstroke$-এর জন্য সেটি নির্গমে দুটি $\TMstroke$
  লিখবে। এগোনোর সঙ্গে সঙ্গে নিবেশ মুছে দিলে নিশ্চিত করা যায় যে কোনো
  $\TMstroke$ বাদ পড়বে না বা দুবার দ্বিগুণ হবে না। গোটা নিবেশ মুছে গেলে
  ফিতায় $2n$টি $\TMstroke$ অবশিষ্ট থাকবে। নির্মিত টুরিং যন্ত্রটির
  দশা-চিত্র \olref{fig:doubler}-এ দেখানো হয়েছে।
  \begin{figure}
\[
\begin{tikzpicture}[->,>=stealth',shorten >=1pt,auto,node distance=2.8cm,
                    semithick]
  \tikzstyle{every state}=[fill=none,draw=black,text=black]

  \node[initial,state] (1)              {$q_0$};
  \node[state]         (2) [right of=1] {$q_1$};
  \node[state]         (3) [right of=2] {$q_2$};
  \node[state]         (4) [below of=3] {$q_3$};
  \node[state]         (5) [left of=4]  {$q_4$};
  \node[state]         (6) [left of=5]  {$q_5$};

  \path (1) edge node {\TMtrans{\TMstroke}{\TMblank}{\TMright}} (2)
    (2) edge [loop above] node {\TMtrans{\TMstroke}{\TMstroke}{\TMright}} (2)
      edge node {\TMtrans{\TMblank}{\TMblank}{\TMright}} (3)
    (3) edge [loop above] node {\TMtrans{\TMstroke}{\TMstroke}{\TMright}} (3)
        edge node {\TMtrans{\TMblank}{\TMstroke}{\TMright}} (4)
    (4) edge [loop below] node {\TMtrans{\TMblank}{\TMstroke}{\TMleft}} (4)
        edge node {\TMtrans{\TMstroke}{\TMstroke}{\TMleft}} (5)
    (5) edge [loop below]  node {\TMtrans{\TMstroke}{\TMstroke}{\TMleft}} (5)
        edge              node {\TMtrans{\TMblank}{\TMblank}{\TMleft}} (6)
    (6) edge [loop below] node {\TMtrans{\TMstroke}{\TMstroke}{\TMleft}} (6)
        edge              node {\TMtrans{\TMblank}{\TMblank}{\TMright}} (1);
\end{tikzpicture}
\]
\caption{একটি দ্বিগুণকারী যন্ত্র}
\ollabel{fig:doubler}
\end{figure}
\end{ex}

\begin{prob}
ইচ্ছামতো একটি নিবেশ বেছে নিয়ে \olref[tur][mac][rep]{ex:doubler}-এর
দ্বিগুণকারী যন্ত্রটির কনফিগারেশনগুলি ধাপে ধাপে অনুসরণ করো।
\end{prob}

\begin{prob}
$\{\TMendtape,\TMblank, A, B\}$ বর্ণমালার এমন একটি টুরিং যন্ত্র নকশা
করো, যা $A$ ও $B$-এর এমন যেকোনো প্রতীকক্রম গ্রহণ করে, অর্থাৎ তাতে থামে,
যেখানে $A$-এর সংখ্যা $B$-এর সংখ্যার সমান \emph{এবং} সব $A$, সব $B$-এর
আগে আসে; আর $A$ ও~$B$-এর সংখ্যা সমান না হলে বা সব~$A$, সব~$B$-এর আগে
না এলে প্রতীকক্রমটি বর্জন করে, অর্থাৎ তাতে থামে না। (যেমন, যন্ত্রটি
$AABB$ ও $AAABBB$ গ্রহণ করবে, কিন্তু $AAB$ ও $AABBAABB$—দুটিই বর্জন
করবে।)
\end{prob}

\begin{prob}
$\{\TMendtape,\TMblank, A, B\}$ বর্ণমালার এমন একটি টুরিং যন্ত্র নকশা
করো, যা নিবেশ হিসেবে $A$ ও $B$-এর যেকোনো প্রতীকক্রম~$\alpha$ নেয় এবং
সেটির অনুলিপি করে $\alpha\alpha$ আকারের নির্গম উৎপন্ন করে। (যেমন,
$ABBA$ নিবেশের ফলে $ABBAABBA$ নির্গম হওয়া উচিত।)
\end{prob}

\begin{prob}
\emph{বর্ণানুক্রমিক?}: $\{\TMendtape,\TMblank, A, B\}$ বর্ণমালার এমন
একটি টুরিং যন্ত্র নকশা করো, যা নিবেশ হিসেবে $A$ ও $B$-এর একটি সসীম
অনুক্রম পেলে সব $A$, সব $B$-এর বাঁ দিকে আছে কি না পরীক্ষা করে। যন্ত্রটি
নিবেশ-প্রতীকক্রমটি ফিতায় রেখে দেবে এবং প্রতীকক্রমটি “বর্ণানুক্রমিক” হলে
থামবে, আর তা না হলে চিরকাল চক্রাকারে চলবে।
\end{prob}

\begin{prob}
\emph{বর্ণানুক্রমে সাজানোর যন্ত্র}: $\{\TMendtape,\TMblank, A, B\}$
বর্ণমালার এমন একটি টুরিং যন্ত্র নকশা করো, যা নিবেশ হিসেবে $A$ ও $B$-এর
একটি সসীম অনুক্রম নিয়ে সেগুলিকে এমনভাবে পুনর্বিন্যস্ত করে, যাতে সব $A$,
সব~$B$-এর বাঁ দিকে থাকে। (যেমন, $BABAA$ অনুক্রমটি $AAABB$ অনুক্রমে এবং
$ABBABB$ অনুক্রমটি $AABBBB$ অনুক্রমে বদলে যাওয়া উচিত।)
\end{prob}

\end{document}
